% Meta: ~400 palabras
\section{Conclusiones}

Se entrega un estimado de rendimiento de cebada para las 197 parcelas del Reto AgroCebada 2026 y un dashboard para consultarlo. El modelo firmado es la mezcla: el promedio de un ensamble (LightGBM, bosque aleatorio y red elástica, 121 variables) y un LightGBM de solo longitud. En la validación estratificada por municipio, que representa a las 59 parcelas ocultas, el RMSE es 0.508 t/ha, el MAE 0.354 y el $R^{2}$ 0.64, frente a 0.846 de predecir el promedio. En zona nueva (LOMO) el RMSE sube a 0.595, con un sesgo cercano a cero, y queda entre el ensamble, que se cae, y la longitud, que generaliza mejor pero explica menos del cultivo.

Esa elección no fue el ganador automático de un barrido amplio. El superciclo y la prueba de datos externos (INEGI y SoilGrids) se evaluaron con un criterio escrito antes de ver el resultado. No bajaron el error estratificado en al menos 0.010 t/ha sin empeorar el LOMO. La longitud sí se incorporó, como la mitad de la mezcla, porque en esta región el gradiente oeste–este es real y medible.

Las preguntas guía tienen respuesta acotada. El índice de mayor contribución SHAP es el NDVI, sobre todo por la senescencia y no por el pico. La variable climática de mayor peso es la lluvia de mayo. El bloque de menor RMSE es el clima, empatado en la práctica con el modelo completo; el satélite solo no basta. Esas lecturas describen un sexto de la mezcla, el LightGBM del ensamble, y la lluvia de mayo puede ser proxy de zona.

Los límites son parte del resultado. Hay una sola temporada, el modelo no debe salir de estos municipios sin datos nuevos, y el simulador atenúa a la mitad cualquier cambio de clima. El producto es una herramienta de apoyo a la consulta y, en el margen, a la conversación de crédito y asistencia técnica. No es un aforo ni una resolución de financiamiento.

El dashboard ya muestra el mapa, la ficha, las métricas y el simulador. Queda abierto repetir el ciclo en otras temporadas, publicar el contraste contra las 59 etiquetas cuando el comité las libere y volver a probar datos de suelo solo si pasan el mismo criterio. El código de la API y del dashboard está publicado, con instrucciones para correrlo en local, de modo que esa revisión no exija reconstruir el trabajo desde cero.
